Label noise lowers the accuracy of linear classifiers.
We measure what a known remedy, an out-of-fold filter, removes, recovers and costs for logistic regression: it drops samples whose given labels disagree with a logistic regression fit on the other folds, then refits once, with no clean validation set, noise-rate estimate or confidence threshold.
On two small datasets, digits and breast cancer, at two symmetric noise rates, most drops are mislabeled, and the kept set is much cleaner than the training set.
In every noisy setting, the refit model's mean test accuracy exceeds that of logistic regression with the same settings (not significantly on breast cancer at the lower rate) but is below the clean-label reference, winning back only part of the lost accuracy.
With clean labels, few samples are removed, at little cost.
